\section{The knot game on PD shadows}
\label{sec:game}

\subsection{PD conventions}
\label{sec:pd}

We use the standard planar diagram convention documented by
KnotInfo~\cite{knotinfo-pd} and the Knot Atlas~\cite{katlas-pd}; see also
Kauffman~\cite{kauffman-mlinks} and Mastin~\cite{mastin2013links} for code
representations, crossing changes, and the relation to Reidemeister moves. A
crossing is written
\[
X[a,b,c,d],
\]
read counterclockwise beginning with the incoming under-strand. Hence
\[
\boxed{(a,c)\text{ is the under-strand}},\qquad
\boxed{(b,d)\text{ is the over-strand}}.
\]
Strand pairing is positional: implementations must never sort the four labels
and reconstruct strands by numerical adjacency. Production inputs are oriented,
one-component PD codes with consecutive one-based edge labels $1,\dots,2c$ for
$c$ crossings, advancing cyclically along the knot ($n \to 1$). Under this
convention a positive crossing has its over-strand oriented $d \to b$
(\texttt{pd\_code\_utils.crossing\_sign}); the cyclic successor rule makes a
naive integer comparison such as \texttt{b > d} incorrect at wraparound
crossings, e.g.\ $[7,1,8,14]$, whose over-strand runs $14 \to 1$. A crossing
change is the one-place cyclic rotation $[a,b,c,d] \to [d,a,b,c]$ for positive
crossings and $[a,b,c,d] \to [b,c,d,a]$ for negative ones
(\texttt{flip\_crossing}), and flipping twice is the identity---a fact proved
in general in the companion Lean formalization and regression-tested in Python.
External PD codes are quotiented by traversal basepoint, component
orientation, and crossing-list order in
\texttt{pd\_code\_utils.canonicalize\_pd\_code} (version
\texttt{oriented-dihedral-v1}), which returns reversible canonical/source
action maps; graph construction pairs the two occurrences of every arc label
directly (serializer \texttt{pd-arc-incidence-v2}) rather than inferring
incidence from tuple slots.

\subsection{Game rules and terminal classification}
\label{sec:rules}

Fix an unresolved shadow with $c$ crossings. A state assigns each crossing one
of three values: unresolved, keep (choice~0), or switch (choice~1), giving
$3^c$ partial states. Action $2i+j$ resolves crossing~$i$ with choice~$j$; play
alternates with the knotter moving first in production
(\texttt{config.knotter\_first = True}). At a terminal state the resolved PD
code is classified by the normalized Jones polynomial computed with an
exhaustive Kauffman-bracket evaluator ($2^c$ smoothing states, cached by PD
code; \texttt{knot\_invariants.py}). The player assigned the unknotter wins
exactly when the polynomial equals~1.

The Jones$=1$ criterion is the \emph{game rule}, not a mathematical claim that
the Jones polynomial detects the unknot in general---it is not known to do so.
Classification is therefore fenced by \texttt{config.max\_validated\_crossings}
(currently seven): raising the bound requires an independent table comparison,
and any out-of-range diagram raises rather than guessing. Terminal evaluation
is fail-loud throughout.

\subsection{Exact solution}
\label{sec:exact}

\texttt{exact\_solver.py} evaluates every state by minimax from the
player-to-move viewpoint ($\pm 1$), with terminal values from the winner
identity above and interior values by negation-maximization over legal actions.
\texttt{evaluate\_exact.py} grades a learned checkpoint against these tables
with the \texttt{SOLVED} criterion: the highest-probability legal action must
be minimax-optimal \emph{and} the value sign must agree with exact minimax at
\emph{every} nonterminal state of each exact-evaluable shadow. Arena win rates
against past champions or random players are diagnostics only; the exact table
is authoritative and the sole basis for any solved claim.

\subsection{Formalization}
\label{sec:lean}

The companion Lean development (\texttt{LeanKnot/}, verified by
\texttt{lake build}) mirrors the Python game logic and certifies its core
combinatorial facts. \texttt{PD.lean} formalizes the oriented cyclic
convention---\texttt{nextLabel}, \texttt{crossingSign} returning an
\texttt{Option} that models the Python \texttt{ValueError} on ill-formed
crossings, and \texttt{flipCrossing}---and proves the general involution
theorem: flipping twice is the identity on label-bounded crossings with a
directed under-strand and a well-defined sign, via a no-2-cycle lemma for
\texttt{nextLabel} when $n > 2$. \texttt{Game.lean} formalizes the
crossing-resolution game (unresolved/keep/switch cells, legal actions
$2i+j$, big-endian terminal indexing matching Python's
\texttt{itertools.product} order) with fuel-bounded minimax and optimal-action
sets. Three shadow games are certified by \texttt{decide}: the trefoil (2 of 8
terminals knotted; knotter-first value $-1$ with all six openings optimal;
unknotter-first value $+1$), the figure-eight (12 of 16 unknot; value $-1$
under both role assignments, hence a second-player win), and $5_2$ (22 of 32
unknot; unknotter wins under both assignments). The stated trust boundary is
explicit: terminal oracle tables are computed by the Python exhaustive
Kauffman-bracket evaluator, and Lean certifies validity, signs, involution,
and minimax values \emph{given} those tables.
